%% Les trois natures d'information : logique, analogique, numérique.
%% Même grandeur physique (courbe en cloche) vue de trois façons.
%%   - logique   : créneau obtenu par comparaison à une valeur seuil ;
%%   - analogique: la grandeur elle-même, continue ;
%%   - numérique : échantillonnage + quantification sur 3 bits (000 à 111).
\documentclass[tikz,border=4pt]{standalone}
\usepackage{liaisons}
\begin{document}
\begin{tikzpicture}[x=1cm,y=1cm,
  axe/.style={-{Stealth[length=4.5pt]}, line width=0.6pt},
  grad/.style={font=\scriptsize},
  titre/.style={font=\small\bfseries, anchor=south}]

%% cloche commune : b(x) = 2 exp(-(x-1.5)^2/0.6), x de 0 à 3
\def\cloche{2*exp(-(\x-1.5)*(\x-1.5)/0.6)}

%% ================= panneau 1 : logique ================================
\begin{scope}[shift={(0,0)}, x=0.82cm, y=1cm]
  \draw[axe] (-0.15,0) -- (3.25,0) node[grad, anchor=north east, inner sep=3pt] {Temps};
  \draw[axe] (0,-0.15) -- (0,2.9);
  % grandeur physique (grise) et seuil
  \draw[black!45, line width=1pt] plot[domain=0:3, samples=60, smooth] (\x,{\cloche});
  \draw[black!70, dash pattern=on 2.5pt off 2pt, line width=0.7pt] (0,1.1) -- (3.1,1.1);
  \node[grad, anchor=east, align=right, text width=0.85cm] at (-0.18,1.1) {valeur seuil};
  % créneau logique
  \draw[solideD, line width=1.4pt, line join=round]
        (0,0.15) -- (0.901,0.15) -- (0.901,2.5) -- (2.099,2.5) -- (2.099,0.15) -- (3,0.15);
  \fill[solideE] (0.901,1.1) circle (2.2pt) (2.099,1.1) circle (2.2pt);
\end{scope}
\node[titre] at (1.23,3.1) {Logique};

%% ================= panneau 2 : analogique =============================
\begin{scope}[shift={(3.4,0)}, x=0.82cm, y=1cm]
  \draw[axe] (-0.15,0) -- (3.25,0) node[grad, anchor=north east, inner sep=3pt] {Temps};
  \draw[axe] (0,-0.15) -- (0,2.9);
  \draw[solideA, line width=1.4pt] plot[domain=0:3, samples=60, smooth] (\x,{\cloche});
\end{scope}
\node[titre] at (4.63,3.1) {Analogique};

%% ================= panneau 3 : numérique ==============================
%% unité verticale = 1 palier de quantification (0,36 cm)
\begin{scope}[shift={(6.9,0)}, x=0.82cm, y=0.36cm]
  \draw[axe] (-0.15,0) -- (3.25,0) node[grad, anchor=north east, inner sep=3pt] {Temps};
  \draw[axe] (0,-0.4) -- (0,8.1);
  \foreach \k/\mot in {0/000,1/001,2/010,3/011,4/100,5/101,6/110,7/111} {
    \draw[line width=0.5pt] (0,\k) -- (-0.1,\k);
    \node[grad, anchor=east, inner sep=2pt] at (-0.12,\k) {\mot}; }
  % grandeur d'origine, en pointillé
  \draw[black!45, line width=0.9pt, dash pattern=on 2pt off 2pt]
        plot[domain=0:3, samples=60, smooth] (\x,{5.5*exp(-(\x-1.5)*(\x-1.5)/0.6)});
  % escalier : valeur échantillonnée puis quantifiée
  \draw[solideA, line width=1.4pt, line join=round]
    (0,0) -- (0.25,0) -- (0.25,1) -- (0.75,1) -- (0.75,2) -- (1,2) -- (1,3) --
    (1.25,3) -- (1.25,4) -- (1.5,4) -- (1.5,5) -- (1.75,5) -- (1.75,4) -- (2,4) --
    (2,3) -- (2.25,3) -- (2.25,2) -- (2.5,2) -- (2.5,1) -- (3,1) -- (3,0);
\end{scope}
\node[titre] at (8.13,3.1) {Numérique};
\end{tikzpicture}
\end{document}
